\chapter{把完成的部分与真正剩下的任务放回一张图}\label{ch:map}
\section{先记住四种不同的结论}
一是必要性：对任何位置数，可达必有 $\rad(G)\mid k$。二是充分性：在指定 $(k,n)$ 上，这个条件也足够。三是最早成功点 $M,H$：还要有覆盖此前维数的反例。四是最终阈值 $N$：需要以后每个维数都成功。

本篇的主要进展集中在第二、三类，同时已经有第四类的统一线性上界。精确最终阈值仍是开放任务。下面的状态按实际定理范围整理；原文件中写于较早阶段的“未解”，只在未被后续结果覆盖的范围内保留。\src{S02}\src{S03}\src{S25}

\section{一般定理的覆盖表}
\begin{longtable}{p{.28\textwidth}p{.37\textwidth}p{.25\textwidth}}
\toprule 范围&已经得到的结论&证据与条件\\\midrule\endhead
全部 $k\ge2$&$M,N\ge B$，$H\ge C$&赋值陷阱与边界同余障碍\\
全部 $k\ge3$&$n\ge4k-1+d_k(n)$ 时G等价；$N\le4k+\lceil\log_2k\rceil+1$&无条件，整数能量下降\\
全部 $k\ge2$&$n=2k+1$ 时G等价&无条件，真实返回与同余群\\
偶数 $k$&全部 $k$ 倍数位置数&无条件，因子转译\\
全部 $k$&全部 $n=ak,a\ge3$&无条件，因子转译\\
奇 $k$，奇 $r<k$，$(k,r)=1$&$n=2k+r$ 时G等价&无条件，奇余量控制器\\
奇素数幂 $k$&全部奇数 $2k<n<3k$&无条件，商块转译\\
$k=t^2$&$n=t^2+t+1$ 时G等价&全参数依赖GRH；小参数、素数模数和充分大参数有无条件结论\\
$k=2s^2$ 或 $3s^2$&$n=k+ms+1$，其中 $m=2,3$&全参数依赖GRH；已证有限小参数无条件\\
$k=ps^2,p\ge5$，$s\ge(p-1)(p-2)$&两种三步可逆反射&资源定理；尚不是完整G充分性\\
全部 $k,n\ge2k+1$&可进入两个 $k$ 块加单点的系统&入口定理；尚不是完整收尾\\
\bottomrule
\end{longtable}

平方族对混合素因子不一定打到候选下界。例如 $k=36$ 时，$S(36)=3$，候选是四十元；平方控制器给四十三元。后者再统一，也没有解决前者。

\section{一些具体元数：已证值与候选值分开看}
表中“已证”均指具体参数的无条件结果；条件性全族的量词另列在上一表。
\begin{longtable}{rrrrp{.39\textwidth}}
\toprule $k$&$B(k)$&$C(k)$&已证 $M,H$&来源或尚缺内容\\\midrule\endhead
2&4&5&$4,5$&前篇二平均定理\\
4&7&7&$7,7$&平方核心 $t=2$\\
6&9&10&$9,10$&四步网络、模五控制器\\
8&13&13&$13,13$&$m=2,s=2$ 的具体控制器\\
9&13&13&$13,13$&平方核心 $t=3$\\
10&16&17&$16,17$&六步网络、模十七控制器\\
12&16&17&$16,17$&四步网络、模十七控制器\\
14&22&22&未定&已有 $M,H\le29$ 的端点界\\
15&21&21&未定&候选尚未覆盖；三十一元端点已解\\
16&21&21&$21,21$&平方族小参数证书\\
18&22&22&未定&$m=2,s=3$ 给二十五元，非二十二元\\
20&25&26&$25,\ ?$&四步网络达到 $M$ 下界\\
25&31&31&$31,31$&平方族小参数\\
27&37&37&$37,37$&$m=3,s=3$ 小参数\\
30&36&37&$36,\ ?$&四步网络达到 $M$ 下界\\
32&41&41&$41,41$&$m=2,s=4$ 小参数\\
36&40&40&未定&$40\le M\le43$，平方点为43\\
42&49&50&$49,\ ?$&四步网络达到 $M$ 下界\\
49&57&57&$57,57$&平方族小参数\\
54&64&65&未定&六十四元G自动合法，仍需可达证明\\
56&64&65&$64,\ ?$&四步网络达到 $M$ 下界\\
64&73&73&$73,73$&平方族小参数\\
72&77&77&未定&$77\le M\le81$，上界来自固定网络\\
81&91&91&$91,91$&平方族小参数\\
120&125&126&未定&一百二十五元G自动合法，仍需固定网络或等价证明\\
128&145&145&$145,145$&$m=2,s=8$ 小参数\\
\bottomrule
\end{longtable}

六平均的旧表需要特别更新。十三元由全部 $2k+1$ 端点覆盖，十四与二十二元由素数因子转译直接覆盖。因此在 $9\le n<24$ 中，当前这些定理还未覆盖的六平均位置数是
\[11,16,17,19,20,23.
\]
这是证明覆盖的空白名单，不是不可解维数名单。仍有 $9\le N(6)\le24$。四平均对应 $7\le N(4)\le16$，其中十一、十三、十五元尚未由上述成果接齐。

还有一条常见混淆：旧 $B$ 公式在 $2\le k\le200$ 中有17个自动G点，分别为
\[
\begin{gathered}
2\to4,\ 6\to9,\ 10\to16,\ 12\to16,\ 20\to25,\ 30\to36,\\
42\to49,\ 54\to64,\ 56\to64,\ 75\to81,\ 90\to96,\\
110\to121,\ 120\to125,\ 132\to144,\ 138\to162,\\
156\to169,\ 182\to196.
\end{gathered}
\]
这张表只计算 $\rad(B(k))\mid k$。其中哪些真有固定网络，要另查构造；算术自动合法没有替代这一证明。

\section{三个主任务的精确增量}
\paragraph{任务一：较大素数的奇次幂。}
目标是在 $k=p^{2a+1},p\ge5$ 的候选位置数
\[n=k+p^{a+1}+1
\]
证明完整G充分性。已有的三步反射覆盖 $a\ge2$ 的一部分资源；需要补完整根与倍率，立方还需补容量可行的返回。可以改变核心或使用状态相关路径，未要求一定延续同一个反射模板。

\paragraph{任务二：混合素因子的真正临界点。}
目标是一般 $M(k)=B(k)$、$H(k)=C(k)$。单个危险赋值已经给共同下界，但不同允许素数的局部要求必须由一条原位置路径同时实现。$k=14,15,18,36$ 的候选点和 $k=54,120$ 的自动点，是不同类型的测试，而不是同一份模群代码换参数。

固定网络可处理自动点，低秩控制器可处理非自动点，因子转译可处理合适整除族。真正需要补的是它们在临界库存内的共同实现规则。复制稳定化表明放宽位置后答案很简单，也指出有限库存是不能删去的条件。

\paragraph{任务三：从成功点延续到所有后续维数。}
目标是 $N(k)=B(k)$，较弱的中间目标是 $N(k)\le2k+1$。目前已证端点、部分奇中间带、整除族和线性尾部；一般合数的偶数余量、多危险素数载体折叠与后续终止仍有空缺。需要一个能连续降低某个良序量的选择规则，或一个能保证存在终端路径的完整结构定理。

\section{去GRH与路径效率是另外两项任务}
平方族的GRH依赖已落在明确的 $H_t$ 生成命题上。更强的无条件小生成元界、利用 $\Phi_3(t)=t^2+t+1$ 的特殊算术，或扩大当前可执行倍率，都可能推进它。已经证明的平方入口和周期无需再搜索。

有限可达性也不等于短路径。固定网络族有显式操作数；大维数整数下降有关于数值大小的界；同余群编译保证有限词，却尚未统一控制关于输入位长的原子词长。现有样本中较短的词可以用于算法改进，不能冒充最短步定理。

\section{后续阅读与接手方式}
主讲义把新增结构按依赖次序重写；来源附录逐条指出原证明、实验与核验入口。附带的原文快照保留技术表、程序链接和历史限定，查一个具体公式时可以直接回到对应来源。

本轮新增的是这份续篇、可追溯的来源包及例子核对。没有宣布新的全体合数定理，也没有重新独立审完全部历史证书。接手者应先选择上述某个精确任务，再看其已有资源与缺口；不必重新从“寻找一个统一矩阵群”开始。
